\documentclass[11pt]{article}
\usepackage[a4paper,margin=22mm]{geometry}
\usepackage{fontspec}
\setmainfont{DejaVu Serif}
\setsansfont{DejaVu Sans}
\usepackage{microtype}
\usepackage{xcolor}
\usepackage{hyperref}
\usepackage{enumitem}
\usepackage{parskip}
\definecolor{Accent}{HTML}{971D32}
\definecolor{Muted}{HTML}{666666}
\hypersetup{colorlinks=true,urlcolor=Accent,linkcolor=Accent}
\setlist{leftmargin=1.6em,itemsep=0.3em,topsep=0.35em}
\setcounter{tocdepth}{2}
\renewcommand{\contentsname}{Contents}
\newcommand{\sectionrule}{\par\vspace{-0.5em}\noindent\textcolor{Accent}{\rule{\linewidth}{0.6pt}}\par}
\begin{document}

\begin{center}
{\sffamily\bfseries\color{Accent} TOEFL WORD OF THE DAY\par}
\vspace{8mm}
{\fontsize{42}{48}\selectfont\bfseries clairvoyant\par}
\vspace{2mm}
{\Large\sffamily\color{Accent} /klɛrˈvɔɪənt/\par}
\vspace{2mm}
{\small\sffamily Local audio phrase: clairvoyant\par}
{\small\sffamily Audio file: audios/clairvoyant.mp3\par}
\vspace{2mm}
{\large\sffamily\color{Muted} adjective and countable noun\par}
\vspace{5mm}
{\sffamily said to perceive hidden or future events \textbullet\ showing unusually accurate foresight \textbullet\ a person said to have such abilities\par}
\vspace{6mm}
{\large Clairvoyant can describe a supposed supernatural ability, an unusually accurate prediction, or a person believed to possess that ability.\par}
\vspace{6mm}
\begin{minipage}{0.84\textwidth}
\itshape “His warning about the economic crisis now seems almost clairvoyant.”
\end{minipage}\par
\vspace{10mm}
\textbf{Date:} July 30th 2026\quad\textbullet\quad \textbf{Level:} B1--mid-B2 TOEFL
\vfill
Created and compiled by \textbf{Amir Kooshky}\par
\vspace{2mm}
\href{https://t.me/KooshkyTOEFL}{Telegram Channel}\quad\textbullet\quad
\href{https://www.instagram.com/kooshkytoefl}{Instagram}\quad\textbullet\quad
\href{https://t.me/KooshkyTOEFL\_pv}{Personal Telegram}
\end{center}
\newpage
\tableofcontents
\newpage

\section{Meanings and usage}
\sectionrule
\subsection{Meaning 1: Able to perceive hidden or future events}
\textbf{Part of speech:} adjective

\textbf{IPA:} /klɛrˈvɔɪənt/

\textbf{Pronunciation phrase:} clairvoyant

\textbf{Local audio file:} audios/clairvoyant.mp3

\textbf{Register:} neutral to formal; often used in discussions of supernatural beliefs, fiction, and psychological research

\textbf{Definition:} believed to be able to know about future, distant, or hidden events without receiving information through the ordinary senses

\textbf{In simpler words:} supposedly able to know hidden things or see the future

\textbf{Typical pattern:} be clairvoyant or have clairvoyant abilities

\subsubsection{Natural examples}
\begin{enumerate}
  \item She claimed to be clairvoyant and able to predict major events.
  \item The story features a clairvoyant woman who sees visions of the future.
  \item I am not clairvoyant, so you need to tell me what you are thinking.
  \item The researchers tested participants who said they had clairvoyant abilities.
  \item According to the legend, the priest was clairvoyant and could see events occurring in distant lands.
  \item He jokingly asked whether she was clairvoyant after her prediction came true.
  \item The novel's main character becomes clairvoyant after a mysterious accident.
\end{enumerate}

\subsubsection{Nuance and implication}
\textbf{Central nuance:} The word refers to knowledge supposedly gained without normal observation, communication, or reasoning.

\textbf{Usual implication:} Writers often use claimed, believed, or supposedly because clairvoyant ability has not been scientifically established.

\textbf{Compared with observant:} An observant person notices information through the ordinary senses. A clairvoyant person is claimed to receive information without those normal sources.

\subsubsection{Grammar notes}
\textbf{How to use it:} Use it after a linking verb such as be or seem, or before a noun such as ability, vision, or character.

\textbf{What can follow it:} It may be followed by enough context to explain the supposed ability, as in clairvoyant about future events, although this pattern is less common than have clairvoyant abilities.

\textbf{Position in a sentence:} The expression I am not clairvoyant is often used humorously to mean that someone cannot know unspoken information or predict what will happen.

\subsubsection{Useful patterns}
\textbf{Pattern:} be clairvoyant

\textbf{Use:} Use this to describe someone who is believed or claimed to have the ability.

\textbf{Example:} The fortune-teller claimed to be clairvoyant.

\textbf{Pattern:} clairvoyant ability

\textbf{Use:} Use this to name the supposed power to perceive hidden or future events.

\textbf{Example:} The study found no reliable evidence of clairvoyant ability.

\textbf{Pattern:} I am not clairvoyant

\textbf{Use:} Use this, often humorously, when someone expects you to know something they have not told you.

\textbf{Example:} I am not clairvoyant; please explain what you need.

\subsubsection{Common wrong usage}
\paragraph{Correction 1}
\textbf{Wrong:} She is a clairvoyant person who can clairvoyance the future.

\textbf{Correct:} She is a clairvoyant person who claims to see the future.

\textbf{Why:} Clairvoyance is a noun, not a verb.

\paragraph{Correction 2}
\textbf{Wrong:} I am not a clairvoyant, so tell me what you want.

\textbf{Correct:} I am not clairvoyant, so tell me what you want.

\textbf{Why:} Use the adjective without a when describing an ability. I am not a clairvoyant is possible only when clairvoyant means a person who practices professionally.

\subsection{Meaning 2: Showing accurate foresight}
\textbf{Part of speech:} adjective

\textbf{IPA:} /klɛrˈvɔɪənt/

\textbf{Pronunciation phrase:} clairvoyant

\textbf{Local audio file:} audios/clairvoyant.mp3

\textbf{Register:} formal and often figurative; common in commentary, reviews, and historical analysis

\textbf{Definition:} showing an unusually accurate understanding of what will happen or of something that is not yet obvious

\textbf{In simpler words:} extremely accurate about what would happen

\textbf{Typical pattern:} seem, appear, or prove clairvoyant

\subsubsection{Natural examples}
\begin{enumerate}
  \item His warning about the economic crisis now seems almost clairvoyant.
  \item The author's clairvoyant analysis anticipated many later technological changes.
  \item Her prediction that remote work would become widespread appeared remarkably clairvoyant.
  \item In retrospect, the committee's concerns about water shortages seem clairvoyant.
  \item The novel offered a nearly clairvoyant picture of life in a highly connected society.
  \item His clairvoyant assessment of the market allowed the company to prepare early.
  \item What once sounded unrealistic now appears to have been a clairvoyant observation.
\end{enumerate}

\subsubsection{Nuance and implication}
\textbf{Central nuance:} This use does not necessarily suggest a real supernatural power. It emphasizes that an earlier prediction or judgment was surprisingly accurate.

\textbf{Usual implication:} The accuracy often becomes clear only later, after the predicted development has occurred.

\textbf{Compared with prescient:} Prescient also means accurately anticipating future events and is usually the more precise formal choice. Clairvoyant is more vivid and may suggest that the accuracy seemed almost supernatural.

\subsubsection{Grammar notes}
\textbf{How to use it:} Use it after seem, appear, or prove, or before nouns such as prediction, warning, analysis, or observation.

\textbf{What can follow it:} The sentence usually explains the later event that made the earlier prediction appear accurate.

\textbf{Position in a sentence:} Modifiers such as almost, nearly, remarkably, and surprisingly are common because this use is usually figurative.

\subsubsection{Useful patterns}
\textbf{Pattern:} seem or appear clairvoyant

\textbf{Use:} Use this when an earlier judgment looks unusually accurate after later events.

\textbf{Example:} Her comments from a decade ago now seem clairvoyant.

\textbf{Pattern:} an almost clairvoyant prediction

\textbf{Use:} Use this for a prediction that was so accurate that it seems supernatural.

\textbf{Example:} The report contained an almost clairvoyant prediction of the energy crisis.

\textbf{Pattern:} a clairvoyant analysis or warning

\textbf{Use:} Use this for an analysis or warning that correctly anticipated later developments.

\textbf{Example:} His clairvoyant warning was ignored by policymakers at the time.

\subsubsection{Common wrong usage}
\paragraph{Correction 1}
\textbf{Wrong:} The prediction was clairvoyance.

\textbf{Correct:} The prediction was clairvoyant.

\textbf{Why:} Use the adjective clairvoyant after was. Clairvoyance is a noun.

\paragraph{Correction 2}
\textbf{Wrong:} Her analysis clairvoyantly the future.

\textbf{Correct:} Her analysis predicted the future with remarkable accuracy.

\textbf{Why:} Clairvoyantly is an adverb and cannot function as the main verb.

\subsection{Meaning 3: Person believed to have supernatural perception}
\textbf{Part of speech:} countable noun

\textbf{IPA:} /klɛrˈvɔɪənt/

\textbf{Pronunciation phrase:} clairvoyant

\textbf{Local audio file:} audios/clairvoyant.mp3

\textbf{Register:} neutral; common in fiction, popular culture, and discussions of supernatural claims

\textbf{Definition:} a person who claims or is believed to know about future, distant, or hidden events through a supernatural ability

\textbf{In simpler words:} a person said to see hidden things or the future

\textbf{Typical pattern:} a clairvoyant who claims to + verb

\subsubsection{Natural examples}
\begin{enumerate}
  \item The clairvoyant claimed that she could describe events occurring far away.
  \item He visited a clairvoyant who offered to predict his future.
  \item The experiment tested whether self-described clairvoyants could identify hidden objects.
  \item In the novel, a young clairvoyant warns the village about an approaching disaster.
  \item The documentary interviewed several people who worked professionally as clairvoyants.
  \item The police rejected the clairvoyant's offer to help locate the missing property.
  \item Skeptics argued that the clairvoyant had obtained the information through ordinary questioning.
\end{enumerate}

\subsubsection{Nuance and implication}
\textbf{Central nuance:} The noun identifies a person associated specifically with perceiving distant, hidden, or future events.

\textbf{Usual implication:} The word itself does not prove that the person's claims are true, so careful writing often includes claimed, alleged, or self-described.

\textbf{Compared with fortune-teller:} A fortune-teller predicts a person's future, often as a profession. A clairvoyant is believed to perceive hidden, distant, or future information and may claim a broader type of supernatural perception.

\subsubsection{Grammar notes}
\textbf{How to use it:} Use a clairvoyant for one person and clairvoyants for more than one.

\textbf{What can follow it:} It is often followed by a clause beginning with who to explain what the person claims to perceive.

\textbf{Noun use:} This is a countable noun, so the singular normally needs a, an, or the.

\subsubsection{Useful patterns}
\textbf{Pattern:} a clairvoyant who claims to + verb

\textbf{Use:} Use this to report the person's claim without presenting it as an established fact.

\textbf{Example:} The program featured a clairvoyant who claimed to communicate with distant people.

\textbf{Pattern:} a self-described clairvoyant

\textbf{Use:} Use this for someone who identifies themselves as having such abilities.

\textbf{Example:} A self-described clairvoyant agreed to participate in the controlled experiment.

\textbf{Pattern:} consult or visit a clairvoyant

\textbf{Use:} Use this when someone seeks advice or predictions from such a person.

\textbf{Example:} The character visits a clairvoyant before beginning her journey.

\subsubsection{Common wrong usage}
\paragraph{Correction 1}
\textbf{Wrong:} She works as clairvoyant.

\textbf{Correct:} She works as a clairvoyant.

\textbf{Why:} The singular countable noun needs a.

\paragraph{Correction 2}
\textbf{Wrong:} The clairvoyance predicted a major change.

\textbf{Correct:} The clairvoyant predicted a major change.

\textbf{Why:} Clairvoyant refers to the person. Clairvoyance refers to the supposed ability.

\section{Word family}
\sectionrule
\subsection{clairvoyance}
\textbf{Part of speech:} uncountable noun

\textbf{IPA:} /klɛrˈvɔɪəns/

\textbf{Definition:} the supposed ability to know about future, distant, or hidden events without using the ordinary senses

\textbf{Relationship to the headword:} This noun names the supposed ability associated with a clairvoyant person.

\textbf{Grammar pattern:} evidence of clairvoyance or claim to possess clairvoyance

\textbf{Usage note:} It is normally used without a or an.

\subsubsection{Examples}
\begin{enumerate}
  \item The researchers found no convincing evidence of clairvoyance.
  \item Clairvoyance plays an important role in the novel's plot.
\end{enumerate}

\subsection{clairvoyantly}
\textbf{Part of speech:} adverb

\textbf{IPA:} /klɛrˈvɔɪəntli/

\textbf{Definition:} in a way that seems to show knowledge of future or hidden events

\textbf{Relationship to the headword:} This adverb describes an action that appears unusually accurate about future or hidden events.

\textbf{Usage note:} It is uncommon. A phrase such as with remarkable foresight is often more natural.

\subsubsection{Examples}
\begin{enumerate}
  \item The author almost clairvoyantly anticipated the rise of digital surveillance.
  \item She seemed to predict the change clairvoyantly.
\end{enumerate}

\section{Common collocations}
\sectionrule
\subsection{clairvoyant ability}
\textbf{Applies to:} Able to perceive hidden or future events

\textbf{Meaning:} the supposed power to perceive future, distant, or hidden events

\textbf{Pattern:} claim, possess, or demonstrate clairvoyant ability

\textbf{Example:} The experiment was designed to test claims of clairvoyant ability.

\subsection{clairvoyant powers}
\textbf{Applies to:} Able to perceive hidden or future events

\textbf{Meaning:} supposed supernatural abilities involving hidden or future knowledge

\textbf{Example:} The character believes that she has clairvoyant powers.

\textbf{Usage note:} The plural powers is common in fictional and popular descriptions.

\subsection{clairvoyant vision}
\textbf{Applies to:} Able to perceive hidden or future events

\textbf{Meaning:} a supposed mental image of a hidden or future event

\textbf{Example:} In the story, a clairvoyant vision warns him of approaching danger.

\subsection{almost clairvoyant}
\textbf{Applies to:} Showing accurate foresight

\textbf{Meaning:} so accurate about the future that it seems nearly supernatural

\textbf{Pattern:} seem or appear almost clairvoyant

\textbf{Example:} Her forecast of the industry's development was almost clairvoyant.

\subsection{clairvoyant prediction}
\textbf{Applies to:} Showing accurate foresight

\textbf{Meaning:} a prediction that proves surprisingly accurate

\textbf{Example:} The essay contained a seemingly clairvoyant prediction about social media.

\subsection{clairvoyant insight}
\textbf{Applies to:} Showing accurate foresight

\textbf{Meaning:} an unusually accurate understanding of a future development

\textbf{Example:} The economist's clairvoyant insight helped explain changes that occurred years later.

\subsection{self-described clairvoyant}
\textbf{Applies to:} Person believed to have supernatural perception

\textbf{Meaning:} a person who describes themselves as having clairvoyant abilities

\textbf{Example:} The study recruited a self-described clairvoyant for a controlled test.

\subsection{consult a clairvoyant}
\textbf{Applies to:} Person believed to have supernatural perception

\textbf{Meaning:} seek advice or information from a person believed to have supernatural perception

\textbf{Example:} In the play, the king consults a clairvoyant before the battle.

\subsection{evidence of clairvoyance}
\textbf{Applies to:} Able to perceive hidden or future events

\textbf{Meaning:} information supporting the existence of the supposed ability

\textbf{Example:} The researchers reported that they had found no reliable evidence of clairvoyance.

\textbf{Usage note:} Clairvoyance is the noun form.

\section{Synonyms and antonyms}
\sectionrule
\subsection{Close synonyms}
\subsubsection{psychic}
\textbf{Part of speech:} adjective

\textbf{Applies to:} Able to perceive hidden or future events

\textbf{Shared meaning:} Both words can describe a person believed to receive information through supernatural means.

\textbf{Difference:} Psychic is broader and can refer to several supposed supernatural mental abilities. Clairvoyant refers more specifically to perceiving distant, hidden, or future events.

\textbf{Example:} She claimed to possess psychic abilities.

\subsubsection{prescient}
\textbf{Part of speech:} adjective

\textbf{Applies to:} Showing accurate foresight

\textbf{Shared meaning:} Both words can describe a prediction or judgment that accurately anticipates later events.

\textbf{Difference:} Prescient means showing accurate knowledge of what will happen and does not usually suggest supernatural powers. Clairvoyant is more vivid and may imply that the accuracy seemed almost supernatural.

\textbf{Example:} The report's prescient warning was largely ignored at the time.

\subsubsection{prophetic}
\textbf{Part of speech:} adjective

\textbf{Applies to:} Showing accurate foresight

\textbf{Shared meaning:} Both can describe words or predictions that later appear unusually accurate.

\textbf{Difference:} Prophetic often refers to a statement that correctly predicts the future and may have a religious or solemn association. Clairvoyant emphasizes apparently seeing what was hidden or still to come.

\textbf{Example:} Her comments about the environmental crisis proved prophetic.

\subsubsection{fortune-teller}
\textbf{Part of speech:} countable noun

\textbf{Applies to:} Person believed to have supernatural perception

\textbf{Shared meaning:} Both words can refer to someone who claims to predict future events.

\textbf{Difference:} A fortune-teller specifically predicts people's futures, often as a service. A clairvoyant may claim to perceive a wider range of hidden, distant, or future events.

\textbf{Example:} The fortune-teller offered to read his cards.

\subsubsection{medium}
\textbf{Part of speech:} countable noun

\textbf{Applies to:} Person believed to have supernatural perception

\textbf{Shared meaning:} Both words refer to people associated with claimed supernatural abilities.

\textbf{Difference:} A medium claims to communicate with the spirits of dead people. A clairvoyant claims to perceive hidden, distant, or future information.

\textbf{Example:} The medium claimed to deliver messages from the dead.

\section{Words learners may confuse}
\sectionrule
\subsection{clairvoyance}
\textbf{Part of speech:} adjective or countable noun versus uncountable noun

\textbf{Why learners confuse them:} The words have closely related meanings and very similar forms.

\textbf{Key difference:} Clairvoyant is an adjective or a noun referring to a person, while clairvoyance is the uncountable noun for the supposed ability.

\textbf{clairvoyant:} The novel describes a clairvoyant child.

\textbf{clairvoyance:} The novel explores the idea of clairvoyance.

\section{Etymology}
\sectionrule
\textbf{Origin:} The word comes from French clairvoyant, formed from words meaning clear and seeing.

\textbf{Early meaning:} Its basic historical sense was clear-sighted or able to see clearly.

\textbf{Development:} It later became associated specifically with the supposed supernatural perception of distant, hidden, or future events.

\section{Practice exercises}
\sectionrule
\subsection{Word family choice}
Choose the best member of the word family for each blank.

\begin{enumerate}
  \item The researchers found no scientific evidence supporting the existence of \_\_\_\_\_.
  \item The novel's main character is a \_\_\_\_\_ who claims to see future events.
  \item His remarkably \_\_\_\_\_ warning anticipated the crisis by several years.
\end{enumerate}

\subsection{Correct or incorrect?}
Decide whether each sentence is correct. Correct the inaccurate ones.

\begin{enumerate}
  \item She claims to possess clairvoyant abilities.
  \item The clairvoyance predicted that the city would expand.
  \item His prediction now seems almost clairvoyant.
  \item She works as clairvoyant in the city center.
\end{enumerate}

\subsection{Collocation practice}
Complete each sentence with a natural collocation.

\begin{enumerate}
  \item The researchers tested her claim of \_\_\_\_\_ ability under controlled conditions.
  \item The author's prediction about digital communication now seems almost \_\_\_\_\_.
  \item The skeptical journalist decided to \_\_\_\_\_ a clairvoyant for the article.
\end{enumerate}

\subsection{Complete answer key}
\subsubsection{Word family choice}
\begin{enumerate}
  \item \textbf{clairvoyance.} The researchers found no scientific evidence supporting the existence of clairvoyance. \emph{Use the noun clairvoyance for the supposed ability.}
  \item \textbf{clairvoyant.} The novel's main character is a clairvoyant who claims to see future events. \emph{Use the countable noun clairvoyant for the person.}
  \item \textbf{clairvoyant.} His remarkably clairvoyant warning anticipated the crisis by several years. \emph{Use the adjective clairvoyant before warning.}
\end{enumerate}

\subsubsection{Correct or incorrect?}
\begin{enumerate}
  \item \textbf{Correct} \emph{Clairvoyant abilities is a natural expression for the supposed supernatural power.}
  \item \textbf{Incorrect}. The clairvoyant predicted that the city would expand. \emph{Use clairvoyant for the person making the prediction.}
  \item \textbf{Correct} \emph{This is a natural figurative use for a prediction that proved unusually accurate.}
  \item \textbf{Incorrect}. She works as a clairvoyant in the city center. \emph{The singular countable noun requires a.}
\end{enumerate}

\subsubsection{Collocation practice}
\begin{enumerate}
  \item \textbf{clairvoyant.} The researchers tested her claim of clairvoyant ability under controlled conditions. \emph{Clairvoyant ability is an established phrase for the supposed power.}
  \item \textbf{clairvoyant.} The author's prediction about digital communication now seems almost clairvoyant. \emph{Almost clairvoyant describes a prediction that proved surprisingly accurate.}
  \item \textbf{consult.} The skeptical journalist decided to consult a clairvoyant for the article. \emph{Consult a clairvoyant means seek information or advice from such a person.}
\end{enumerate}

\vfill
\begin{center}
\small\color{Muted} Created and compiled by \textbf{Amir Kooshky}\\
\href{https://t.me/KooshkyTOEFL}{Telegram Channel}\quad\textbullet\quad
\href{https://www.instagram.com/kooshkytoefl}{Instagram}\quad\textbullet\quad
\href{https://t.me/KooshkyTOEFL\_pv}{Personal Telegram}
\end{center}
\end{document}
